\subsection{划分}

定义6. 若集合A与B不相交（或不交），则称它们为不相交的（或不相交的）；若 \(\mathbf{A} \cap \mathbf{B} = \phi\) 。如果 \(\mathbf{A} \cap \mathbf{B} \neq \phi\) ，则称A与B相遇（或相交）。设 \((\mathbf{X}_{t})_{i \in I}\) 设为一个集合族。该族中的集合被称为互不相交的，如果满足条件 \(i \in I\) , \(x \in I\) , \(i \neq x\) 蕴含 \(\mathbf{X}_{t} \cap \mathbf{X}_{x} = \phi\) .

设 \(f\) 是从 \(\mathbf{A}\) 到 \(\mathbf{B}\) 的映射，并设 \((\mathbf{Y}_i)_{i\in \mathbf{I}}\) 是 \(\mathbf{B}\) 的一族两两不相交的子集。命题4表明， \((\overline{f}^1\langle \mathbf{Y}_i\rangle)_{i\in \mathbf{I}}\) （ \(\mathbf{A}\) 的子集族）中的各集合两两不相交。另一方面，如果 \((\mathbf{X}_i)_{i\in \mathbf{I}}\) 是 \(\mathbf{A}\) 的一族两两不相交的子集，则族 \((f\langle \mathbf{X}_i\rangle)_{i\in \mathbf{I}}\) 中的各集合一般并不两两不相交。

命题8. 设 \((\mathbf{X}_t)_{t \in I}\) 是一族互不相交的集合，且设 \((f_t)_{t \in I}\) 是定义域为 \(f_t\) 是 \(\mathbf{X}_t\) 对于每一个 \(t \in I\) 那么存在恰好一个函数 \(f\) ，其定义域为 \(\bigcup_{i \in I} \mathbf{X}_t\) ，目标为 F，并且它扩展了每一个函数 \(f_t\) ( \(t \in I\) ).

这是命题7的一个推论，因为 \(f_{\iota}\) 并且 \(f_{x}\) 显然在 \(\mathbf{X}_{\iota} \cap \mathbf{X}_{x} = \emptyset\) 每当 \(\iota \neq x\) .

定义7. 集合E的一个划分是E的非空互不相交子集的一个族，它们覆盖E。

例. 对于每个非空集合A，族 \((\{x\})_{x\in\mathbf{A}}\) 是A的一个划分。

如果 \((\mathbf{X}_i)_{i\in I}\) 是集合的一个划分 \(E\) ，映射 \(\iota \to X_i\) 的 \(I\) 到集合 \(\mathfrak{F}\) 的元素 \(X_i\) 的划分是双射的。因此，若 \(\mathfrak{F}\) 给定，则该划分在指标集之间的一一对应意义下是确定的。通常当我们谈到一个划分时，我们所关心的是划分中元素所构成的集合。

